% JAC1964-20261121. /root/arithmetic_sign.
\subsection*{AR-15a. The algebraic radical and a precise grade-zero obstruction}

\begin{proposition}[AR-15a]
Let \(J\), \(N\), and \(\mathcal A_F\) be as in AR-15 for the Fermat
threefold. Suppose that, for every Jacobian \(B\) over a finite extension
of \(F\), the ordinary cycle-class map is injective on
\[
             \operatorname{CH}^2_{(0)}(B\times J)_{\mathbf Q}.
             \tag{AR15a.1}
\]
Then \(\mathcal A_F=0\). More precisely, any nonzero algebraically
trivial class in \(\operatorname{CH}^2(N_F)\) would produce, over a
finite extension, a nonzero homologically trivial class in
\(\operatorname{CH}^2_{(0)}(B\times J)_{\mathbf Q}\) for some Jacobian
\(B\).
\end{proposition}

\begin{proof}
Use the curve-family expression in the proof of AR-15. If its sum is
nonzero, at least one projected summand
\(\Gamma_*D\), with \(D\in\operatorname{Pic}^0(C)_{\mathbf Q}\),
is nonzero. Insert the degree-one Chow projector of \(C\), since it
acts identically on \(D\), and the projector of \(N\) on the target.
The resulting correspondence is nonzero because its action on \(D\)
is nonzero.

Let \(B=\operatorname{Jac}(C)\) and \(b=\dim B\). The actual isomorphism
\(h^1(C)\simeq h^1(B)\) transports this correspondence to a nonzero
\[
 \Gamma_B\in\operatorname{CH}^{b+1}(B\times J)_{\mathbf Q},
 \qquad
 \Gamma_B=\pi_N\circ\Gamma_B\circ\pi_1^B.
 \tag{AR15a.2}
\]
Here the codimension is \(b+1\), because the induced cohomological
map raises degree one on \(B\) to degree three on \(J\).
Its cohomological realization is zero: a nonzero map
\(H^1(B)\to H(N)(1)\) would have a nonzero image of types
\((1,0),(0,1)\), while \(H(N)(1)\) has types
\((2,-1),(-1,2)\).

As a cycle on the product, (AR15a.2) is supported in the Chow--K\"unneth
component
\[
        \pi_{2b-1}^B\otimes\pi_3^J.
\]
The first index is \(2b-1\), not one, because the source projector
is transposed when a correspondence is viewed as a cycle on the
product. Thus the total index is \(2b+2=2(b+1)\).
Equivalently \(\Gamma_B\) belongs to Beauville grade zero.

The Fourier isomorphism on the abelian variety \(B\times J\), of
dimension \(b+3\), sends a grade-zero class of codimension \(b+1\)
to one of codimension
\[
             (b+3)-(b+1)=2.
\]
It preserves homological triviality and is invertible on rational
Chow groups. Therefore the Fourier transform of \(\Gamma_B\) is a
nonzero homologically trivial class in the space in (AR15a.1), as
claimed. Under the displayed injectivity hypothesis this is
impossible, proving the proposition.
\end{proof}

\paragraph{Status and implication.}
The Fourier isomorphism and the degree-one Chow-motive isomorphism
are the same established inputs used in AR-1a and AR-2.
Injectivity in (AR15a.1) is stated here as an additional hypothesis;
it has not been deduced from CM Hodge algebraicity, numerical
positivity, Kimura nilpotence, or the new Selmer bound.
This proposition gives a concrete obstruction to constructing a
nonzero algebraically trivial zero-height direction in \(N\):
such a direction would also violate grade-zero cycle-class
injectivity on an explicit abelian product. It does not prove
that the remaining Griffiths quotient is generated by the
positive CM orbit, or that its real height is positive.
